% Terjemahan Bahasa Indonesia dari prelude.tex baris 185-205.
% Sumber: Methods of Algebra, Volume 2, commit
% 9a5803ff77dd3257484cb177f851a73770a59dd3 (CC BY 4.0).
% stable-unit-id: o014.aljabr2.prelude.acknowledgments
% Catatan penerjemah: Romanisasi sebelas nama dalam daftar ucapan terima kasih
% dan nama editor sengaja ditangguhkan hingga ejaan Latinnya dapat diverifikasi
% dari sumber primer; nama-nama itu dipertahankan persis dalam aksara Han pada
% teks pembaca.

% segment-id: o014.li2.prelude.acknowledgments.h001
\section*{Ucapan Terima Kasih}
\addcontentsline{toc}{section}{Ucapan Terima Kasih}
\hypertarget{o014.aljabr2.prelude.acknowledgments}{ }

% segment-id: o014.li2.prelude.acknowledgments.p001
Selama penyusunan buku ini, penulis menerima koreksi dan saran dari para guru,
rekan, dan sahabat. Karena prosesnya berlangsung lama, kekeliruan ingatan tak
terelakkan. Berikut dicantumkan sebuah daftar yang tidak lengkap: 郇真, 黄晨欣,
黎景辉, 雷嘉乐, 李长远, 孙毓泽, 汤一鸣, 薛寒玉, 阳恩林, 张好风, 赵泽龙
(diurutkan menurut Hanyu Pinyin).

% segment-id: o014.li2.prelude.acknowledgments.p002
Sebagian kecil isi buku ini pernah digunakan dalam perkuliahan di University
of Chinese Academy of Sciences dan Huazhong University of Science and
Technology. Penulis juga menyampaikan terima kasih kepada semua peserta
perkuliahan tersebut.

% segment-id: o014.li2.prelude.acknowledgments.p003
Seperti semua karya penulis hingga kini, buku ini ditulis sepenuhnya dalam
lingkungan perangkat lunak sumber terbuka, termasuk sistem operasi dan rupa
huruf yang digunakan sebelum naik cetak. Penulis menyampaikan penghormatan
kepada begitu banyak pengembang tanpa pamrih; mereka mewujudkan nilai-nilai
universal tentang berbagi dan altruisme.

% segment-id: o014.li2.prelude.acknowledgments.p004
Selama penulisan buku ini, penelitian penulis didukung oleh proyek
\emph{Excellent Young Scientists Fund} 11922101 dari National Natural Science
Foundation of China, dengan Peking University sebagai institusi tuan rumah.
Untuk itu penulis juga menyampaikan terima kasih.

% segment-id: o014.li2.prelude.acknowledgments.p005
Penulis juga berterima kasih kepada editor 赵天夫 dari Higher Education Press
atas kesabaran dan profesionalismenya. Penulisan buku ini semestinya
menyambung Jilid I tanpa terputus, tetapi akibat berbagai hambatan baru dapat
diselesaikan pada penghujung 2022.

% segment-id: o014.li2.prelude.acknowledgments.p006
Meskipun pengerjaannya tertunda, seluruh penyusunan buku ini menguras lebih
banyak tenaga batin penulis daripada Jilid I, terutama pemeriksaan dan
revisi berulang-ulang pada tahap akhir, seolah berjalan di terowongan gelap
tanpa ujung. Sebagian yang cukup besar dari pekerjaan tahap akhir ini dilakukan
di dalam gerbong kereta bawah tanah Beijing. Bekerja di kereta jauh dari nyaman dan
menuntut berpacu dengan waktu dalam arti harfiah. Namun, di dalam arus besar
tanpa nama yang berulang hari demi hari ini, penulis seakan merasakan sesuatu
yang jarang dijumpai di dunia kerja: sebentuk saling pengertian dalam diam,
dan menemukan setitik penghiburan dalam keremangan. Apakah demikian, atau
bukan? Penjelasannya mungkin mengada-ada, tetapi perasaan itu tetaplah nyata.
Preludium akan berakhir; hormat kepada kejauhan tanpa batas dan insan-insan
tanpa akhir.

% segment-id: o014.li2.prelude.acknowledgments.sig001
\vspace{1em}
\begin{flushright}\begin{minipage}{0.3 \textwidth}
	\begin{tabular}{c}
		% segment-id: o014.li2.prelude.acknowledgments.sig002
		{\fangsong Wen-Wei Li} \\
		% segment-id: o014.li2.prelude.acknowledgments.sig003
		Desember 2022 \\
		% segment-id: o014.li2.prelude.acknowledgments.sig004
		di Shijingshan
	\end{tabular}
\end{minipage}\end{flushright}
